\chapter{Introduction}
\label{chap:introduction}

\chapterabstract{This chapter motivates zero-shot static 3D human-scene interaction synthesis from language, a calibrated RGB image, and a reconstructed scene mesh. It places the thesis in relation to learned HSI models, scene-constrained body optimization, and recent foundation-model-based systems, then states the problem formulation, approach, and contributions.}

\section{Why Human-Scene Interaction Matters}

Indoor scenes are built around human action. Chairs are placed so that people can sit, tables so that people can work or eat, cabinets so that objects can be stored and retrieved, and handles so that drawers, doors, and appliances can be opened. The meaning of these objects is therefore not exhausted by their geometry or category labels. A chair is not only a shape occupying space; it affords support for a seated body. A drawer is not only a box-like structure inside a cabinet; it has a front face and a handle that invite a particular hand motion. A room becomes meaningful because human bodies can stand, reach, grasp, lean, sit, lie, and move through it.

This distinction is central to 3D scene understanding. Modern reconstruction methods can produce increasingly accurate models of indoor environments, including camera poses, dense geometry, object surfaces, and scene layout. However, a reconstructed scene by itself remains a passive description of the environment. It tells us where surfaces are, but not which surfaces can support the body. It shows object boundaries, but not which regions are functional for a task. It captures the state of a room, but not the ways in which a person could physically use that room. For applications that require human-centered reasoning, this missing layer is fundamental.

Human-scene interaction provides this missing layer by connecting scene structure to bodily action. Instead of asking only what objects are present, it asks how a person can engage with them. Where should the feet be placed when reaching for a cabinet? Which part of a sofa should support the pelvis and back during sitting? Where should the hand make contact when opening a drawer? What body orientation is plausible when lifting a trash bin? These questions require reasoning jointly about object function, human body parts, contact, balance, reachability, and spatial constraints. They cannot be answered by object recognition or scene reconstruction alone.

This capability is important for a broad class of downstream systems. In embodied AI and robotics, an agent must reason not only about object locations, but also about how humans act around those objects and which parts of the scene are functionally relevant. In augmented and virtual reality, inserted humans must appear grounded in the environment rather than floating near objects or intersecting with them. In animation, synthetic data generation, and digital twins, realistic human activity requires bodies that are placed according to the affordances of the scene. In human-centered scene analysis, the presence or absence of plausible interactions can reveal whether a reconstructed environment is usable, accessible, or semantically coherent.

The challenge is therefore not simply to add a human body to a 3D scene. A visually plausible body can still be physically or functionally wrong: a seated person may hover above a chair, a hand may reach the side of a cabinet instead of the drawer handle, or a body may intersect with furniture while appearing reasonable from one camera view. Human-scene interaction requires a stronger form of consistency. The body must be semantically aligned with the requested action, geometrically aligned with the scene, and physically plausible with respect to contact and collision.

This thesis studies human-scene interaction synthesis as a step toward scene representations that are not only geometric, but actionable. The goal is to synthesize human bodies that express how a real indoor scene can be used. By grounding human pose in scene affordances, object function, and physical contact, HSI synthesis moves beyond static reconstruction toward a more complete understanding of indoor environments as spaces for human activity.

\section{From Reconstructed Geometry to Interaction Synthesis}

A 3D reconstruction converts visual observations of a room into a metric representation of space. It provides camera poses, scene geometry, object surfaces, and the spatial layout of the environment. This representation is essential for grounding perception in real-world coordinates: it tells us where the floor is, how far a chair is from a table, where a cabinet is located, and which surfaces are visible or occluded. However, reconstruction alone does not specify how a human body should be arranged in the scene. Geometry describes the shape of the environment, but interaction requires reasoning about the relationship between that geometry and the body.

This gap becomes apparent even for simple activities. To synthesize a person sitting on a chair, it is not enough to know that a chair exists. The body must be placed so that the pelvis and thighs are supported by the seat, the feet remain near the floor, and the torso has a plausible orientation relative to the chair. To synthesize a person opening a drawer, it is not enough to locate the cabinet. The hand should reach the drawer pull or handle, while the body should stand at a reachable distance and face the object in a plausible way. To synthesize a person lifting a trash bin, the hands, arms, torso, and feet must form a coherent configuration around the object. In each case, the interaction is defined not only by the object category, but by the spatial and functional relation between body parts and scene parts.

Human-scene interaction synthesis addresses this problem by generating a human body configuration that is jointly consistent with an action and a 3D scene. Unlike image-based human insertion, the goal is not only to make a rendered view appear plausible. The synthesized body must exist in the same coordinate frame as the scene and must satisfy constraints that are meaningful in 3D: contact should occur at appropriate locations, the body should have a plausible scale and orientation, and severe interpenetration with the environment should be avoided. The output is therefore a scene-grounded human state, not simply a visual overlay.

This distinction separates interaction synthesis from several related tasks. Object detection identifies what is present in the scene, but it does not determine how the object can be used. Human pose estimation recovers a body configuration from observations, but it usually assumes that the person is already present. Text-to-image generation can produce visually convincing depictions of people interacting with objects, but the result is not necessarily metrically consistent with a reconstructed scene. Human-scene interaction synthesis combines aspects of all three problems while adding a stronger requirement: the generated body must be semantically aligned with the requested action and physically grounded in the 3D environment.

The difficulty is that interaction is structured. A single instruction often implies a set of relations among the action, the object, the relevant object part, the participating body parts, and the supporting scene surfaces. Sitting involves support relations between the lower body and a seat. Opening involves a reaching or grasping relation between a hand and a functional element. Lifting involves coordinated contacts between the hands and the object, together with a stable stance on the floor. These relations are not directly given by the reconstruction. They must be inferred from the task and then realized as a plausible body placement in the scene.

Interaction synthesis can therefore be viewed as the step from a passive scene model to an actionable scene model. A passive model answers geometric questions such as where surfaces are and how objects are arranged. An actionable model supports questions about what a human can do in the environment, where the body should be placed, which body parts should contact which scene regions, and whether the resulting configuration is physically credible. This thesis focuses on this transition: synthesizing a human body that does not merely appear inside a reconstructed room, but uses the room in a way that is consistent with language, geometry, and physical contact.

\section{The Zero-Shot Challenge}

Human-scene interaction synthesis becomes especially difficult when the system is expected to operate in a zero-shot setting. In a closed setting, the model may assume a fixed set of object categories, a fixed set of actions, or training examples that closely resemble the test scenes. Under these assumptions, the problem can be approached as learning regularities from paired examples of humans interacting with objects. The model can learn, for instance, that people sit on chairs, stand on floors, touch tables with their hands, or lie on beds. However, such regularities are only reliable when the test interaction lies within the distribution of the training data.

Real indoor scenes do not satisfy this assumption. They contain diverse object instances, cluttered layouts, partial visibility, unusual viewpoints, and functional elements that vary widely in appearance. A drawer handle may be horizontal, vertical, recessed, metallic, wooden, large, small, or partially occluded. A chair may be an office chair, dining chair, stool, sofa, or bench. Even when the object category is familiar, the region relevant to an interaction may differ across instances. The same action phrase can therefore require different geometric solutions depending on the scene. `Open the drawer'' requires the hand to reach the handle of a particular drawer in a particular cabinet. `Sit on the chair'' requires the body to find a support surface of the correct height and orientation. ``Lift the bin'' requires a graspable region on the visible object and a stable full-body configuration around it.

This variability exposes the limitation of relying only on supervised 3D human-scene interaction data. Paired HSI data is expensive to collect because it requires reconstructed scenes, accurate human body fits, object annotations, action labels, contact information, and enough diversity in body shapes, object instances, and room layouts. As a result, existing datasets cover only a small subset of the possible interactions that occur in indoor environments. Learned HSI models can produce strong results for common interactions represented in their training sets, but they often depend on closed vocabularies, predefined object categories, known target objects, or interaction patterns observed during training. These constraints make it difficult to synthesize interactions for arbitrary language instructions in previously unseen reconstructed scenes.

Recent work has therefore explored the use of large vision-language, image-generation, and video-generation models as sources of interaction knowledge. These models have seen broad visual and linguistic evidence of people using objects, and they can often suggest plausible body poses, object affordances, and interaction layouts without being trained directly on 3D HSI datasets. This has motivated zero-shot and training-free approaches to human-scene interaction synthesis, where foundation models provide semantic or visual priors and 3D geometry provides metric grounding. Methods such as GenZI, FunHSI, and ZeroHSI show that this direction is promising because interaction knowledge can be transferred from large-scale visual models to 3D scenes without requiring paired 3D interaction supervision.

However, zero-shot HSI is not solved by generating a plausible image of a person interacting with an object. A generated image may express the intended action visually, but it does not guarantee that the body is correctly placed in the reconstructed scene. The hands may appear near an object in image space while missing the corresponding 3D surface. The body may look plausible from one camera view while penetrating furniture or floating above the floor. The generated interaction may also contact the wrong part of the object: the side of a cabinet instead of the drawer handle, the rim of a bin instead of a graspable side, or the back of a chair instead of the seat. The challenge is therefore to convert broad interaction priors into scene-specific, metrically valid human-body configurations.

The zero-shot setting thus requires solving several coupled questions. Given an open-vocabulary instruction, the system must infer which scene object is relevant, which part of that object affords the action, which human body parts should participate, where supporting contacts should occur, and how the full body should be arranged in 3D. These questions cannot be answered independently. The choice of contact region affects the pose; the pose affects the feasible body location; the body location affects reachability, balance, and collision; and the scene geometry constrains all of these decisions. Zero-shot HSI is therefore not simply a problem of text-conditioned pose generation, but one of grounding an inferred interaction structure in a specific reconstructed environment.

This thesis addresses the zero-shot challenge in the static 3D setting. Rather than assuming a fixed action vocabulary or paired 3D interaction examples for every task, it treats the instruction as an open-ended description of how a person should use the scene. The goal is to infer the relevant interaction structure and synthesize a single SMPL-X body that realizes that structure in the original scene coordinate frame. This formulation keeps the problem narrower than full dynamic interaction generation, but it preserves the central difficulty of zero-shot HSI: the system must transform language and scene observations into a physically plausible human-scene contact configuration without relying on task-specific 3D supervision.

\section{Body-Part--Scene-Part Grounding}

The previous sections framed human-scene interaction synthesis as more than placing a plausible body near a relevant object. The remaining question is what structure should connect the language instruction, the human body, and the scene. A natural first answer is to reason at the object level: a person sits on a chair, opens a drawer, touches a table, or lifts a bin. Object-level reasoning is useful, but it is not precise enough for interaction synthesis. Most interactions are not defined by whole objects. They are defined by relations between specific parts of the human body and specific regions of the scene.

Consider the instruction `a person opening a drawer.'' The target object may be a cabinet or drawer unit, but the interaction itself is concentrated at the drawer handle or front edge. The relevant human part is usually a hand, not the entire body. At the same time, the rest of the body is not arbitrary: the feet must remain supported by the floor, the torso must be oriented toward the cabinet, and the arm must be arranged so that the hand can reach the functional part. Similarly, `a person sitting on a chair'' primarily requires support between the lower body and the seat, but also implies a plausible relation between the feet and the floor and between the torso and the chair orientation. These examples show that an interaction is best understood as a set of body-part--scene-part relations.

This thesis adopts that view. Instead of treating the human body as a single object and the scene as a collection of object categories, it represents an interaction through participating body regions, scene regions, and contact relations between them. The human body is decomposed into interaction-relevant parts, such as hands, arms, feet, pelvis, thighs, back, and torso. These parts correspond to regions on the SMPL-X surface that commonly participate in physical interactions. On the scene side, the relevant regions may be broad support surfaces, such as floors, seats, beds, and tabletops, or smaller functional elements, such as handles, knobs, switches, faucets, and drawer pulls. The task of HSI synthesis is then to arrange the body so that the appropriate body regions are grounded to the appropriate scene regions.

This formulation is important because different actions require different kinds of contact. Support interactions, such as sitting, lying, standing, or leaning, depend on stable contact between load-bearing body regions and supporting scene surfaces. Manipulation interactions, such as opening, pulling, pressing, or lifting, depend on localized contact between the hands and functional object parts. Many interactions combine both types. A person opening a drawer must stand on the floor while reaching the handle. A person lifting a bin must grasp the object while maintaining a stable stance. A person leaning over a sink may contact the floor with the feet, the counter with the hands or torso, and the faucet with one hand. Representing these interactions as structured body-scene contacts makes these dependencies explicit.

The body-side representation also matters. SMPL-X provides a detailed parametric human mesh, but the mesh vertices themselves are not semantically meaningful unless they are grouped into interaction regions. This thesis therefore uses manually defined contact regions on the SMPL-X body surface corresponding to interaction parts such as hands, arms, feet, pelvis, back, and other task-relevant regions. These regions provide a bridge between high-level language and low-level geometry: an instruction may mention a hand, a seated posture, a leaning action, or a grasping motion, while the final synthesis must place specific surface regions of the body near specific surface regions of the scene.

The scene-side representation is equally important. For many tasks, detecting the object category or even the object part is insufficient. The same cabinet can support several different interactions: opening a drawer, reaching for an upper shelf, leaning on the counter, or standing beside it. Each task selects a different scene region and a different set of body regions. Conversely, the same body part can participate in different functional roles: a hand can grasp a handle, press a switch, support the body on a tabletop, or hold the side of a bin. Interaction synthesis must therefore resolve not only which object is involved, but which object region is relevant and what type of body contact it requires.

Moreover, part-level grounding alone does not fully determine the interaction. Many object parts are spatially extended, and the exact location of contact within a detected part can strongly affect the plausibility of the synthesized pose. For example, when lifting a table, it is not enough to identify the tabletop as the relevant part; the hands should contact plausible grasp or support locations on the tabletop, often near accessible edges or corners rather than at arbitrary points on the surface. Similarly, for an interaction involving a ladder, the system must reason not only about the ladder as an object, or about its rails and rungs as parts, but also about which specific rail and rung regions should be contacted by the hands and feet. The same rung may support a foot at one location while a hand grasps a nearby rail at another, and these local choices determine whether the resulting pose is physically and functionally coherent.

This distinction is important because semantic part labels provide structure, but they do not by themselves specify contact geometry. A tabletop, rail, rung, seat, counter, cabinet front, or floor region can contain many possible contact locations, only some of which are compatible with the requested action, the human pose, the object affordance, and the surrounding scene. Body-part--scene-part grounding therefore requires a finer notion of localization: the framework must identify not only the relevant object and object part, but also the appropriate contact region within that part. This finer grounding allows the synthesized body to interact with the scene in a way that is both semantically correct and geometrically plausible.

To capture this structure, the thesis uses a Scene Interaction Graph as an intermediate representation. The graph expresses the target scene element, the participating human interaction parts, and the expected contact or support relations between them. Its purpose is not to describe every detail of the final pose, but to specify the interaction constraints that the pose should satisfy. In this sense, the graph acts as a compact bridge between open-language task descriptions and 3D body synthesis. It turns an instruction such as ``lift the trash bin with both hands'' into a structured interaction hypothesis involving the bin, the two hands, the arms, the feet, and the supporting floor.

A key aspect of this thesis is that the scene-side contact regions are not assumed to be given by manual annotations, object geometry alone, or part labels alone. Instead, they are inferred through image-based generative reasoning. Starting from the scene image, a diffusion-based inpainting model first synthesizes a plausible human performing the requested interaction in the scene. This generated human provides a visual hypothesis about pose, orientation, laterality, and approximate contact layout. Then, conditioned on the interaction parts specified by the Scene Interaction Graph, a second diffusion-based inpainting step marks the corresponding contact regions on the target object or support surface in the original scene image. In this way, the scene image is used twice: first to imagine how the interaction should look, and then to localize where the relevant body parts should make contact on the object or scene.

This image-based contact grounding is particularly valuable for functional interactions, where the relevant contact region may be small, localized, or ambiguous even when the correct object part is known. A cabinet may contain many possible surfaces, but only the handle or drawer pull is appropriate for opening. A trash bin may have several visible faces, but only some are plausible grasp regions for lifting. Likewise, a large tabletop, ladder rail, or ladder rung may be semantically correct at the part level, but the interaction still depends on where within that part the hands, feet, or other body regions should make contact. By using diffusion models to propose both the human interaction hypothesis and the object-side contact regions, the framework leverages broad visual priors while keeping the final interaction tied to the actual scene and to the fine-grained spatial layout of the interaction.

Body-part--scene-part grounding is therefore the central abstraction behind this thesis. It makes the interaction explicit enough to guide synthesis, but flexible enough to support open-vocabulary instructions and previously unseen scenes. By grounding actions in relations between human surface regions and scene surface regions, and by further localizing contact within those scene regions, the formulation moves beyond object-level and part-level semantics toward a representation of how the body physically and functionally engages with the environment.

\section{Problem Setting and Contributions}

This thesis studies zero-shot static 3D human-scene interaction synthesis in real reconstructed indoor scenes. The input consists of a natural-language instruction, a calibrated RGB image of the scene, and a reconstructed scene mesh. The output is a single SMPL-X body placed in the original scene coordinate frame. The synthesized body should express the requested action, contact the appropriate scene regions with the appropriate body parts, and remain physically plausible with respect to scale, orientation, support, and collision.

The setting is zero-shot because the method is not trained on paired 3D examples of the requested interaction in the target scene. The instruction may describe common support interactions, such as sitting on a chair, lying on a bed, or leaning on a table, or functional interactions, such as opening a drawer, pressing a switch, turning a faucet, or lifting a bin. In each case, the system must infer the relevant scene element, the functional or supporting region on that element, the participating human body parts, and the spatial arrangement required to realize the interaction. The target object and the required contact regions are not assumed to be available as ground-truth annotations.

The formulation is static. The goal is not to synthesize a full motion sequence, simulate forces, or modify the state of the scene. Instead, the output represents a single interaction state. For steady actions such as sitting or lying, this state corresponds to the interaction itself. For actions involving movable or articulated objects, such as opening a drawer, pulling a door, or lifting a trash bin, the synthesized body represents the pre-action or contact-ready state: the human has reached the relevant scene region and adopted a plausible pose, but the object has not yet moved. This choice isolates a fundamental subproblem of human-scene interaction: determining where contact should occur and how the body should be configured in a real 3D environment.

The proposed approach is built around body-part--scene-part grounding. An open-language instruction is first interpreted as a structured interaction involving human interaction regions, scene regions, and contact relations. On the human side, SMPL-X is treated not only as a parametric body model, but as a surface with semantically meaningful interaction regions, such as hands, arms, feet, pelvis, thighs, back, and torso. On the scene side, the relevant regions may be broad support surfaces or localized functional elements. The synthesis problem is then to place the body so that these body regions and scene regions are brought into physically and semantically meaningful correspondence.

To support this grounding in open-vocabulary scenes, the thesis combines structured interaction reasoning with generative visual priors and scene-aware 3D fitting. A Scene Interaction Graph represents the expected relations among the action, target scene element, participating body parts, and contact regions. Diffusion-based image reasoning provides visual hypotheses about how the interaction should appear and where the corresponding object-side contacts should lie. These hypotheses are then grounded in the reconstructed scene to synthesize a metric SMPL-X body that satisfies the inferred interaction structure. At a high level, the method uses foundation models to reason about what interaction is plausible, while using the scene reconstruction to ensure that the final body is placed in 3D rather than merely depicted in an image.

This thesis makes the following contributions:

\begin{enumerate}
\item \textbf{A zero-shot formulation of static 3D human-scene interaction synthesis.}
The thesis defines the task of synthesizing a single SMPL-X body in a real reconstructed indoor scene from an open-language instruction, a calibrated RGB image, and a scene mesh, without requiring paired 3D interaction examples for the target task.

\item \textbf{A body-part--scene-part grounding view of HSI.}
The thesis formulates interaction synthesis as the problem of grounding relations between human interaction regions and scene interaction regions, rather than only matching whole-body poses to object categories.

\item \textbf{A Scene Interaction Graph representation for open-vocabulary interactions.}
The thesis introduces a structured representation that captures the target scene element, participating human body parts, scene regions, and expected contact or support relations implied by a natural-language instruction.

\item \textbf{A SMPL-X interaction-part abstraction for contact-aware synthesis.}
The thesis defines semantically meaningful contact regions on the SMPL-X body surface, enabling high-level interaction descriptions involving hands, arms, feet, pelvis, back, and other body regions to be connected to low-level 3D body geometry.

\item \textbf{An image-guided strategy for localizing scene-side contact regions.}
The thesis uses generative image reasoning to infer where the human interaction parts should contact the target object or support surface, allowing functional regions such as handles, knobs, drawer pulls, and graspable object areas to be identified without manual scene-side contact annotation.

\item \textbf{A scene-grounded synthesis framework for physically plausible body placement.}
The thesis synthesizes a full SMPL-X body in the reconstructed scene coordinate frame, using the inferred interaction structure and scene geometry to encourage semantic consistency, appropriate contact, realistic scale, and reduced physical artifacts such as floating or penetration.

\end{enumerate}

Together, these contributions move static HSI synthesis beyond object-level placement and toward interaction-level grounding. The thesis argues that plausible human-scene interaction requires knowing not only what object is involved, but which parts of the human body should engage which parts of the scene, and how that engagement can be realized as a coherent 3D body configuration.

\section{Thesis Structure}

\newcommand{\structurechapter}[2]{%
\par\medskip\noindent\hyperref[#1]{\textbf{Chapter~\ref*{#1}}} #2\par
}

\structurechapter{chap:related-work}{reviews the foundations and related work necessary for this thesis, including 3D scene reconstruction, parametric human body models, human-scene interaction synthesis, affordance and contact reasoning, foundation models for visual generation, and evaluation of physically grounded human-scene configurations.}

\structurechapter{chap:methodology}{presents the proposed zero-shot static HSI synthesis framework. It describes the Scene Interaction Graph, the SMPL-X interaction-part abstraction, diffusion-based human and contact-region reasoning, scene-side contact grounding, and the final scene-aware SMPL-X fitting procedure.}

\structurechapter{chap:experiments}{defines the experimental setup, including the ScanNet++ scenes, interaction instructions, selected action categories, implementation details, baselines or ablations, and quantitative and qualitative evaluation metrics.}

\structurechapter{chap:results}{reports the experimental results and analyzes the behavior of the method. It examines semantic consistency, contact accuracy, physical plausibility, pose diversity, qualitative examples, ablation studies, and common failure modes.}

\structurechapter{chap:conclusion}{summarizes the findings of the thesis and discusses limitations and future directions, including dynamic interaction synthesis, richer physical reasoning, multi-view grounding, multi-human interactions, and more complete modeling of object state changes.}
